\chapter{The Editors}\label{cha:editors}

Two of them, because they answer different questions. \verb|EDIT| is the
one to reach for when a file wants a line changed and you want to be out
again in ten seconds. \verb|VI| is the one for a file too big to think
about, and it is modal, so it expects you to have met \texttt{vi} before.

Both draw through JIM, the VT100 in hardware (Chapter~\ref{cha:tube}), and
both take their keys raw from the ROM --- an arrow arrives as one byte
rather than an escape sequence to unpick. Neither needs the Tube, so both
run on the Pi as well as the desktop.

\section{EDIT}
\verb|EDIT name| opens a file, or starts an empty one if there is no such
file yet. The whole text sits in memory below the program, so it holds
about 22\,KB --- ample for a \texttt{STARTUP.BAT}, a \texttt{.SUB}, a
BASIC listing or a letter.

\begin{type}
arrows Home End PgUp PgDn   move
Backspace Delete            rub out
Enter                       split the line
Ctrl-O  or  Ctrl-S          save
Ctrl-X                      leave
\end{type}

The bar along the bottom carries the name, a \verb|*| while there are
unsaved changes, where the cursor is, and those last two keys, so there is
nothing to remember. Backspace at the start of a line joins it to the one
above. Long lines slide sideways rather than wrap. \verb|Ctrl-X| on a
changed file says so and waits: press it again to throw the changes away,
or \verb|Ctrl-O| to keep them.

\section{VI}
\verb|VI name| is modal in the old way: the keys are commands until
\verb|i| (or \verb|a|, \verb|I|, \verb|A|, \verb|o|, \verb|O|) puts you
into insert, and Escape brings you back out.

\begin{type}
counts   3dd 5j 2dw 10G      before almost anything
move     h j k l arrows  w b e  0 ^ $  gg G  PgUp PgDn
operate  d c y + any motion, or doubled: dd cc yy
insert   i a I A  o O  s S  C            Esc leaves insert
change   x X r ~ J D  p P (the unnamed register)
undo     u        redo Ctrl-R
search   /pat  ?pat  n N
ex       :w  :q  :q!  :wq  :x
         :s/old/new/[g]    :%s/old/new/[g]
         :map lhs rhs      :imap lhs rhs
\end{type}

A count goes in front of nearly anything: \verb|3dd| deletes three lines,
\verb|2dw| two words, \verb|10G| goes to line ten. Operators take any
motion --- \verb|d| \verb|c| \verb|y| with \verb|w|, \verb|b|,
\verb|e|, \verb|0|, \verb|^|, \verb|$|, \verb|G| --- or double
themselves to work on whole lines. What \verb|d|, \verb|x| or \verb|y|
took goes into the unnamed register for \verb|p| and \verb|P|. Undo is
unlimited, and so is redo: the journal lives in far memory with the file,
so one level would have cost the same as all of them.

\verb|:map| and \verb|:imap| define your own keys --- \verb|:imap jk|
\verb|<Esc>| is the usual one --- and last for the session.

Two behaviours are deliberate and worth knowing before they surprise you.
A charwise operator \textbf{clamps to the line}: \verb|dw| at the end of
a line stops there instead of eating into the next one. And a search is a
\textbf{plain substring}, not a pattern --- \verb|/foo| finds
\verb|foo|, and \verb|.| \verb|*| \verb|[| mean themselves. So does
the left-hand side of \verb|:s|.

\verb|:q| on a changed file refuses and says so; \verb|:q!| throws the
changes away and \verb|:wq| keeps them. \verb|:w name| writes somewhere
else. Lines past the end of the file are marked \verb|~| down the left, as
they should be.

\begin{aside}
\textbf{Where the file lives.} Nothing of it is in the 64\,KB. Every line
is a 256-byte slot in far memory --- a length and up to 255 characters ---
and only the line under the cursor is brought down, fetched when the
cursor arrives and written back when it leaves. So the ceiling is far
memory rather than the address space: 32000 lines, and \verb|LOAD| and
\verb|SAVE| go straight to and from a flat copy out there without the file
ever passing through the CPU's view.

That is the whole design, and it is deliberately the \emph{only} design:
there is no small-file case kept separately to disagree with the large one
at exactly the wrong moment. It costs nothing because the DMA engine
copies with \texttt{memmove} semantics, so opening or closing a line is a
single transfer of everything below it however long the file is.

A 1.28\,MB file of 20000 lines --- twenty times the whole address space
--- opens, edits at the far end and saves back byte for byte.
\end{aside}

\section{Editing without losing what you were doing}
A program loaded from the shell lands on top of whatever was in memory, so
running an editor from inside a BASIC would ordinarily destroy the
program you were writing. \verb|SWAP| (Chapter~\ref{cha:shell}) is the way
round it:

\begin{type}
*SWAP EDIT MYPROG.BAS
\end{type}

from EhBASIC puts the machine away, runs the editor on a clean one, and
gives yours back --- program, variables, screen and all --- when the
editor exits.
